\setchapterpreamble[u]{\margintoc}
\chapter{Linear Transformations and Change of Basis}
\labch{linear-transformations}

\sources{Strang, \emph{Introduction to Linear Algebra}, 6th ed.\ \cite{strang2023ila},
§8.1--8.3; MIT 18.06 \cite{ocw1806}, Lectures~33--34.}

A matrix is a table of numbers; a linear transformation is a rule. This chapter
separates the two, which explains why the same rule can be recorded by many
different matrices and why choosing the basis well is worth the trouble.

\section{What makes a map linear}
\labsec{linear-maps}

\begin{definition}
A map \(T:V\to W\) between vector spaces is \emph{linear} when
\[
	T(\vect{u}+\vect{v})=T(\vect{u})+T(\vect{v}),
	\qquad
	T(c\vect{v})=cT(\vect{v})
\]
for all \(\vect{u},\vect{v}\in V\) and scalars \(c\).
\end{definition}

Setting \(c=0\) gives \(T(\vect{0})=\vect{0}\), so any map that shifts the origin
is disqualified. Rotation about the origin, projection onto a subspace, and
differentiation of polynomials are linear; translation and the map
\(\vect{v}\mapsto\lVert\vect{v}\rVert\) are not.

\marginnote{The two conditions combine into one: \(T\) is linear exactly when it
commutes with taking linear combinations. Everything in this book is a
consequence of that single permission.}

\section{The matrix of a transformation}
\labsec{transformation-matrix}

Fix a basis \(\vect{v}_1,\ldots,\vect{v}_n\) of \(V\) and a basis
\(\vect{w}_1,\ldots,\vect{w}_m\) of \(W\). Because \(T\) respects combinations,
it is determined by what it does to the basis vectors. Writing each image in the
output basis,
\[
	T(\vect{v}_j)=\sum_{i=1}^{m}a_{ij}\vect{w}_i,
\]
collects the coefficients into a matrix \(A\in\R^{m\times n}\), and the rule
becomes matrix multiplication in coordinates.

\begin{remark}
Column \(j\) of \(A\) is nothing other than the coordinate vector of
\(T(\vect{v}_j)\). Building the matrix of a transformation is always the same
mechanical act: feed in each basis vector, write the answer in the output basis,
and stack the results as columns.
\end{remark}

\begin{example}
Let \(V\) be the polynomials of degree at most \(3\), with basis
\(1,t,t^{2},t^{3}\), and let \(W\) be those of degree at most \(2\), with basis
\(1,t,t^{2}\). Differentiation sends \(1\mapsto0\), \(t\mapsto1\),
\(t^{2}\mapsto2t\), \(t^{3}\mapsto3t^{2}\), so
\[
	A=\begin{bmatrix}0&1&0&0\\0&0&2&0\\0&0&0&3\end{bmatrix}.
\]
Its nullspace is spanned by \([1,0,0,0]\tp\), the constants --- exactly the
polynomials that differentiate to zero --- and its rank is \(3\), matching
\(3+1=4\) in \refch{four-subspaces}.
\end{example}

\begin{example}
Rotation of \(\R^{2}\) by an angle \(\theta\) sends \(\vect{e}_1\) to
\((\cos\theta,\sin\theta)\) and \(\vect{e}_2\) to \((-\sin\theta,\cos\theta)\), so
\[
	A=\begin{bmatrix}\cos\theta&-\sin\theta\\ \sin\theta&\cos\theta\end{bmatrix},
\]
which satisfies \(A\tp A=I\): rotations preserve length, as an orthogonal matrix
must.
\end{example}

\section{Change of basis}
\labsec{change-of-basis}

Keeping the transformation and changing the basis changes the matrix. Let \(M\)
have as its columns the new basis vectors expressed in the old basis. A vector
with new coordinates \(\vect{c}\) has old coordinates \(M\vect{c}\), and the two
matrices of the same transformation are related by
\[
	A_{\text{new}}=M^{-1}A_{\text{old}}M .
\]

This is the similarity relation of \refsec{similar}, now with its meaning
supplied: similar matrices are the same transformation described by two
observers using different rulers. Eigenvalues, trace, determinant and rank are
properties of the transformation and are therefore identical for all of them;
the individual entries are not.

\marginnote{A useful sanity check when a formula looks wrong: ask whether the
quantity you computed could possibly depend on the choice of basis. If it
should not, and your expression does, the error is in the bookkeeping.}

\section{Choosing the basis well}
\labsec{good-basis}

Diagonalisation is the statement that, for a diagonalisable \(A\), the
eigenvector basis makes the matrix diagonal:
\(\Lambda=X^{-1}AX\). In that basis the transformation is \(n\) independent
scalings, and questions about powers, exponentials, and long-run behaviour become
questions about numbers.

When \(A\) is not square, or not diagonalisable, the singular value
decomposition supplies the substitute. Taking \(\vect{v}_1,\ldots,\vect{v}_n\) as
the input basis and \(\vect{u}_1,\ldots,\vect{u}_m\) as the output basis, the
matrix of the transformation is \(\Sigma\) --- diagonal, non-negative, and
ordered. No two-basis choice does better.

\begin{remark}
Read this way, the factorisations of Part~I are all answers to one question:
\emph{which coordinates make this transformation simple?} Elimination answers it
for solving, Gram--Schmidt for projecting, diagonalisation for iterating, and
the singular value decomposition for approximating.
\end{remark}
